% =====================================================================
%  AEGIS — one-page summary (ARGUS Open Innovation Challenge 2026)
%  Compile with XeLaTeX:   xelatex AEGIS_2.tex
%  All fonts are bundled in ./fonts (Fira Sans family + Font Awesome 7
%  Free Solid), so nothing needs to be installed system-wide.
%
%  Layout note: every box and text block is placed at absolute page
%  coordinates (unit = 1bp, origin = top-left corner of the A4 page),
%  taken from the original PDF, so the output lines up 1:1 with it.
%  Paragraphs are still normal TeX paragraphs (justification, Greek
%  hyphenation and microtype protrusion are done by TeX), so you can edit
%  the text; just keep in mind that boxes do not grow by themselves.
% =====================================================================
\documentclass[10pt]{article}
\usepackage[a4paper,margin=0pt]{geometry}
\usepackage{fontspec}
\usepackage{polyglossia}
\setmainlanguage{greek}
\setotherlanguage{english}
\usepackage{xcolor}
\usepackage{tikz}
\usepackage{microtype}
\usepackage{eso-pic}
\usepackage{etoolbox}

% ---------------------------------------------------------------- fonts
\defaultfontfeatures{Path=fonts/, Ligatures=TeX}
\setmainfont{FiraSans-Regular.ttf}[BoldFont=FiraSans-SemiBold.ttf]
\newfontfamily\lightfont{FiraSans-Light.ttf}
\newfontfamily\condsemi{FiraSansExtraCondensed-SemiBold.ttf}
\newfontfamily\condbold{FiraSansExtraCondensed-Bold.ttf}
\newfontfamily\condxbold{FiraSansExtraCondensed-ExtraBold.ttf}
\newfontfamily\monofont{FiraMono-Regular.ttf}
\newfontfamily\FA{FontAwesome7Free-Solid-900.otf}
% Font Awesome icons by code point
\newcommand\fa[1]{{\FA\symbol{"#1}}}

% --------------------------------------------------------------- colors
\definecolor{aegis}{HTML}{0B5D4B}     % main green
\definecolor{aegisdark}{HTML}{084839} % header bottom strip
\definecolor{aegisline}{HTML}{14705C} % header decoration lines
\definecolor{ink}{HTML}{14211B}       % dark text / badge rings
\definecolor{muted}{HTML}{5F6E66}     % small caps labels
\definecolor{soft}{HTML}{EDF2EB}      % light green panels
\definecolor{frame}{HTML}{CFD9D1}     % card border
\definecolor{paper}{HTML}{FBFCFA}     % page background

\pagestyle{empty}
\setlength\parindent{0pt}
\setlength\parskip{0pt}

% --------------------------------------------------------------- helpers
% \T(x,y){text}      : text whose baseline origin sits at (x,y)
% \Par(x,y){w}{text} : justified paragraph, first baseline at (x,y)
% \Rag(x,y){w}{text} : ragged-right paragraph, first baseline at (x,y)
% Text is collected while the drawing is built and placed on top of it at
% shipout with \put, at exact positions (bp -> sp with exact arithmetic).
\def\pagetext{}
\def\spx#1{\fpeval{round((#1)*65781.76)}}
\def\spy#1{\fpeval{round(-(#1)*65781.76)}}
\def\T(#1,#2)#3{\gappto\pagetext{\put(\spx{#1},\spy{#2}){#3}}}
\def\Par(#1,#2)#3#4{\gappto\pagetext{\put(\spx{#1},\spy{#2})%
  {\parbox[t]{#3bp}{\leavevmode#4}}}}
\def\Rag(#1,#2)#3#4{\gappto\pagetext{\put(\spx{#1},\spy{#2})%
  {\parbox[t]{#3bp}{\raggedright\leavevmode#4}}}}
% rounded box: \RBox{color}{radius}(x0,y0)(x1,y1)
\def\RBox#1#2(#3,#4)(#5,#6){\fill[#1,rounded corners=#2bp] (#3,#4) rectangle (#5,#6);}
% section badge (dark ring / white ring / green disc) with a white icon
\def\Badge(#1,#2)#3{%
  \fill[ink]   (#1,#2) circle (8.787bp);
  \fill[white] (#1,#2) circle (7.937bp);
  \fill[aegis] (#1,#2) circle (6.662bp);}
% big stat badge
\def\BigBadge(#1,#2){%
  \fill[ink]   (#1,#2) circle (11.055bp);
  \fill[white] (#1,#2) circle (10.205bp);
  \fill[aegis] (#1,#2) circle (8.929bp);}
% font-size shortcuts
\def\bodyA{\fontsize{9.7}{12.6}\selectfont}    % problem box
\def\bodyB{\fontsize{9.2}{11.9}\selectfont}    % cards
\def\bodyC{\fontsize{8.4}{10.6}\selectfont}    % step cards
\def\headfont{\condbold\fontsize{13.5}{16}\selectfont\color{ink}}

% shortcuts used in the stat tiles, step cards and the features line
\def\statnum{\condbold\fontsize{15}{18}\selectfont\color{ink}}
\def\statlbl{\monofont\fontsize{6.1}{7}\selectfont\color{muted}}
\def\staticon{\fontsize{9}{10}\selectfont\color{white}}
\def\stepnum{\condbold\fontsize{10}{12}\selectfont\color{white}}
\def\steptitle{\condbold\fontsize{11}{13}\selectfont\color{ink}}
\def\stepicon{\fontsize{9.5}{11}\selectfont\color{aegis}}
\def\chev(#1){\T(#1,541.520){\fontsize{8}{9}\selectfont\color{aegis}\fa{232A}}}
\def\ico#1{{\color{aegis}\fa{#1}}\,}

\begin{document}
\null
\begin{tikzpicture}[remember picture,overlay,x=1bp,y=-1bp,
                    shift={(current page.north west)}]

% ================================================================ page
\fill[paper] (0,0) rectangle (595.28,841.89);

% ============================================================== header
\fill[aegis] (0,-0.01) rectangle (595.28,93.54);
\begin{scope}
  \clip (0,0) rectangle (595.28,93.54);
  \foreach \r in {40,52,64,76,88,100}
    \draw[aegisline,line width=0.35pt] (657.64,172.91) circle (\r mm);
  \foreach \a in {95,110,125,140,155,170}
    \draw[aegisline,line width=0.35pt] (657.64,172.91) -- ++(\a:110mm);
\end{scope}
\fill[aegisdark] (0,90.703) rectangle (595.28,93.538);

% shield + temple logo
\draw[white,line width=4.2169bp,line join=round]
  (41.993,21.072) .. controls (52.907,21.072) and (61.835,18.096) .. (66.796,14.127)
  .. controls (71.757,18.096) and (80.686,21.072) .. (91.600,21.072)
  -- (91.600,43.891)
  .. controls (91.600,58.773) and (80.686,68.694) .. (66.796,73.655)
  .. controls (52.907,68.694) and (41.993,58.773) .. (41.993,43.891)
  -- cycle;
\fill[white] (51.914,33.969) -- (66.796,26.529) -- (81.679,33.969) -- cycle;
\fill[white] (51.914,34.962) rectangle (81.679,36.946);
\foreach \x in {53.899,59.603,65.308,71.013,76.718}
  \fill[white] (\x,38.435) rectangle ++(2.976,14.386);
\fill[white] (50.922,54.308) rectangle (82.671,56.292);

\T(102.247,44.782){\condxbold\fontsize{36}{40}\selectfont\color{white}AEGIS}
\T(103.381,69.727){\lightfont\fontsize{13.2}{16}\selectfont\color{white}Ο επισκέπτης γίνεται φύλακας του μνημείου}

% pills
\RBox{white}{3.118}(353.909,22.381)(393.655,38.269)
\T(360.712,33.864){\condsemi\fontsize{9.6}{11}\selectfont\color{aegis}Track A}
\draw[white,line width=0.6pt,rounded corners=3.118bp] (398.121,21.200) rectangle (555.096,39.450);
\T(404.925,33.103){\condsemi\fontsize{9.6}{11}\selectfont\color{white}ARGUS Open Innovation Challenge 2026}

% footer bar
\fill[aegis] (0,829.984) rectangle (595.28,841.89);

% ========================================================== stat tiles

\RBox{soft}{4.535}(39.685,106.299)(162.285,145.060)
\BigBadge(60.378,123.874)
\T(55.895,127.254){\staticon\fa{1F3DB}}
\T(77.953,123.670){\statnum 5}
\T(77.9529695964,136.943){\statlbl ΜΝΗΜΕΙΑ ΤΗΣ ΕΥΡΩΠΗΣ}

\RBox{soft}{4.535}(170.787,106.299)(293.387,145.060)
\BigBadge(191.480,123.874)
\T(186.997,127.236){\staticon\fa{F51E}}
\T(208.756,123.741){\statnum 12–16}
\T(209.055,137.014){\statlbl ΝΟΜΙΣΜΑΤΑ}

\RBox{soft}{4.535}(301.890,106.299)(424.490,145.060)
\BigBadge(322.583,124.298)
\T(318.660,127.941){\staticon\fa{F073}}
\T(340.1579543947,125.267){\statnum έως 7}
\T(340.1579695964,138.540){\statlbl ΜΕΡΕΣ ΩΣ ΦΥΛΑΚΑΣ}

\RBox{soft}{4.535}(432.992,106.299)(555.592,145.060)
\BigBadge(453.685,123.874)
\T(448.641,127.201){\staticon\fa{F1AB}}
\T(471.260,123.670){\statnum 5}
\T(471.260,136.943){\statlbl ΓΛΩΣΣΕΣ}

% ========================================================= the problem
\Badge(48.472,165.753){}
\T(44.885,168.667){\fontsize{7.2}{8}\selectfont\color{white}\fa{26A0}}
\T(62.6459543947,169.375){\headfont Το πρόβλημα}

\RBox{white}{4.535}(39.685,177.657)(555.597,257.598)
\draw[aegis,line width=3.6851bp] (41.528,177.657) -- (41.528,257.598);
\Par(56.693,195.095){485.87}{\bodyA\color{ink}%
Κανείς δεν αμφισβητεί ότι το ARGUS έχει κάνει μια σπουδαία δουλειά.
Μετρά, υπολογίζει και ειδοποιεί για τους κινδύνους που απειλούν πέντε
ευρωπαϊκά μνημεία. Η προστασία όμως, δεν εξαρτάται μόνο από τους
αισθητήρες. \textbf{Η προστασία των μνημείων εξαρτάται και από όσους τα
επισκέπτονται.} Το ARGUS δεν έχει κάποιο μέσο ευαισθητοποίησης των
πολιτών, οπότε οι επισκέπτες των μνημείων δεν γνωρίζουν τι πρέπει να
προσέξουν, ή γιατί υπάρχουν κανόνες για τους επισκέπτες. Χωρίς την
ευαισθητοποίηση του κοινού, η διάσωσή τους μένει ημιτελής.}

% ===================================================== why it matters
\RBox{frame}{6.367}(39.685,266.846)(291.972,390.832)
\RBox{white}{5.669}(40.382,267.543)(291.274,390.135)
\Badge(62.209,286.536){}
\T(59.519,289.226){\fontsize{7.2}{8}\selectfont\color{white}\fa{F252}}
\T(76.3829239911,290.159){\headfont Γιατί υπάρχει}
\Par(53.422,306.835){224.815}{\bodyB\color{black}%
Η φθορά ενός μνημείου είναι αργή και αθόρυβη. Η υγρασία ανεβαίνει στους
τοίχους για χρόνια, το αλάτι βγαίνει στην πέτρα, το ψηφιδωτό χάνει
ψηφίδες. Όταν φανεί με γυμνό μάτι, η ζημιά έχει ήδη γίνει. Η γνώση που
την προλαβαίνει μένει στους ειδικούς. Έτσι το κοινό βλέπει μόνο το
αποτέλεσμα, ποτέ τη φροντίδα, και δεν καταλαβαίνει ότι μετράει και το
δικό του χέρι.}

% ============================================================ our idea
\RBox{aegis}{5.669}(303.307,266.846)(555.594,390.832)
\fill[white] (325.135,285.839) circle (8.788bp);
\T(322.445,288.417){\fontsize{7.2}{8}\selectfont\color{aegis}\fa{1F4A1}}
\T(339.3079239911,289.461){\condbold\fontsize{13.5}{16}\selectfont\color{white}Η ιδέα μας}
\Par(316.347,306.137){226.207}{\bodyB\color{white}%
Σκεπτόμενοι πώς μπορούμε να εξηγήσουμε τη δουλειά του ARGUS,
συνειδητοποιήσαμε πως τα λόγια δεν ωφελούν. Γι’ αυτό, μέσω του
\textbf{AEGIS}, δώσαμε σε οποιονδήποτε την ευκαιρία να \textbf{γίνει
φύλακας} στα μνημεία του ARGUS για έως και 7 μέρες, δίνοντας ρεπό σε
έναν επινοημένο φύλακα του μνημείου.}

% ======================================================== our approach
\RBox{soft}{5.669}(39.685,400.078)(555.597,657.692)
\Badge(61.512,418.506){}
\T(57.925,421.196){\fontsize{7.2}{8}\selectfont\color{white}\fa{F4D7}}
\T(75.6859391929,422.129){\headfont Πώς το προσεγγίσαμε}
\Par(52.725,438.804){489.83}{\bodyB\color{black}%
Η προσέγγισή μας βασίστηκε στα δεδομένα που μας δόθηκαν. Οι κίνδυνοι, οι
θέσεις των αισθητήρων και οι συναγερμοί προέρχονται από το ARGUS.}

% step cards: \Step{card x0}{number}{nx}{title}{icon}{icon x}{text}

\RBox{white}{4.535}(52.725,458.279)(165.404,618.783)
\fill[aegis] (69.166,474.155) circle (7.370bp);
\T(67.218,477.517){\stepnum 1}
\T(81.0719543947,477.777){\steptitle Τοποθέτηση}
\T(146.867,477.777){\stepicon\fa{F51E}}
\Rag(61.796,491.739){94.537}{\bodyC\color{black}%
Ο παίκτης, με 12 έως 16 νομίσματα, ανάλογα με το επίπεδο κινδύνου,
καλείται να επιλέξει ανάμεσα σε μια γκάμα αισθητήρων και να τους
τοποθετήσει πάνω στο εκάστοτε μνημείο.}
\chev(169.716)

\RBox{white}{4.535}(178.442,458.279)(291.121,618.783)
\fill[aegis] (194.883,474.155) circle (7.370bp);
\T(192.786,477.592){\stepnum 2}
\T(206.788,477.778){\steptitle Συναγερμός}
\T(273.767,477.778){\stepicon\fa{1F514}}
\Rag(187.5127719733,491.740){94.537}{\bodyC\color{black}%
Μόλις ολοκληρώσει την τοποθέτηση, έρχεται αντιμέτωπος με έναν πραγματικό
συναγερμό του ARGUS και εξετάζει αν το δίκτυό του θα είχε ανιχνεύσει τον
κίνδυνο. Αμέσως μετά, του εμφανίζονται οι αισθητήρες που τοποθέτησε το
ARGUS, πού και γιατί.}
\chev(295.433)

\RBox{white}{4.535}(304.158,458.279)(416.837,618.783)
\fill[aegis] (320.599,474.154) circle (7.370bp);
\T(318.432,477.506){\stepnum 3}
\T(332.5048175786,477.777){\steptitle Η ζημιά σε 3D}
\T(398.300,477.777){\stepicon\fa{F1B2}}
\Rag(313.229,491.739){94.537}{\bodyC\color{black}%
Τέλος, του εμφανίζεται ως εικονογράφηση η ζημιά που θα γινόταν, πάνω σε
πραγματικό 3D μοντέλο του ARGUS, αν κανείς δεν το πρόσεχε.}
\chev(421.1494831394)

\RBox{white}{4.535}(429.875,458.279)(542.554,618.783)
\fill[aegis] (446.316,474.155) circle (7.370bp);
\T(443.955,477.607){\stepnum 4}
\T(458.221,477.778){\steptitle Ο φύλακας}
\T(522.834,477.778){\stepicon\fa{F505}}
\Rag(438.9458327804,491.740){94.537}{\bodyC\color{black}%
Ο φύλακας, που χάρη στον παίκτη πήρε ρεπό, επιστρέφει με εκπαιδευτικό
ύφος και του εξηγεί γιατί είναι σημαντικό, όταν επισκεπτόμαστε ένα
μνημείο, να ακολουθούμε τους κανόνες που έχουν οριστεί.}

% features line
\Rag(52.7252584303,633.039){489.83}{\bodyB\color{black}%
Όλα τα παραπάνω είναι διαθέσιμα σε \ico{1F310}\textbf{5 γλώσσες}, τα
κείμενα \ico{1F50A}\textbf{διαβάζονται φωναχτά}, το παιχνίδι παίζεται σε
\ico{1F4BB}\textbf{υπολογιστή} και \ico{F3CF}\textbf{κινητό} και είναι
κατάλληλο για τη \ico{1F3EB}\textbf{σχολική τάξη}.}

% =================================================== what the user gains
\RBox{ink}{6.367}(39.685,666.939)(291.972,802.973)
\RBox{white}{5.669}(40.382,667.636)(291.274,802.276)
\Badge(62.209,686.630){}
\T(59.071,689.264){\fontsize{7.2}{8}\selectfont\color{white}\fa{1F464}}
\T(76.3828479822,690.252){\headfont Τι κερδίζει ο χρήστης}
\Par(53.422,706.928){224.815}{\bodyB\color{black}%
Καταλαβαίνει εις βάθος τις ανάγκες και τους κινδύνους ενός μνημείου.
Επιπλέον, φεύγει με έναν \textbf{κώδικα επισκέπτη} και μια \textbf{λίστα}
με όσα να προσέχει στην επόμενη επίσκεψη. Αν δει κάτι, το γράφει σε μια
\textbf{κάρτα παρατήρησης} και τη στέλνει, αν θέλει, στον φορέα του
μνημείου. Ο εκπαιδευτικός βλέπει τις επιλογές της τάξης δίπλα σε όσα
εγκατέστησε το ARGUS.}

% ================================================== what ARGUS gains
\RBox{ink}{6.367}(303.307,666.939)(555.594,802.973)
\RBox{white}{5.669}(304.004,667.636)(554.896,802.276)
\fill[ink]   (325.832,686.630) circle (8.788bp);
\fill[white] (325.832,686.630) circle (7.937bp);
\fill[aegis] (325.832,686.630) circle (6.661bp);
\T(322.245,689.320){\fontsize{7.2}{8}\selectfont\color{white}\fa{1F4E1}}
\T(340.0048935875,690.253){\headfont Τι κερδίζει το ARGUS}
\Par(317.044,707.131){224.813}{\bodyB\color{black}%
Αποκτά ένα \textbf{δημόσιο πρόσωπο} για τα δεδομένα του και πολίτες
εκπαιδευμένους σε αυτά. Οι παρατηρήσεις των πολιτών αποθηκεύονται ήδη
στη δική του μορφή εγγραφής πολιτιστικής τεκμηρίωσης, έτοιμες αν το
ARGUS αποφασίσει να τις δεχτεί. Το παιχνίδι φέρνει κάθε μέρα τους
νεότερους ελέγχους κινδύνου του ARGUS, οπότε ο παίκτης βλέπει την πιο
πρόσφατη κατάσταση του μνημείου.}

\end{tikzpicture}%
\AddToShipoutPictureFG*{\AtPageUpperLeft{\setlength\unitlength{1sp}\selectlanguage{greek}\pagetext}}
\end{document}
